\section{Dataset}

The starting corpus contains 37,554 raw news articles (from \texttt{data/raw/jsons}).
Triplets are constructed in \texttt{data/sort\_data.py} by embedding each article,
retrieving nearest neighbors, and then using an anchor-based enumeration strategy:
for each anchor article $i$, all valid neighbor pairs $(j, k)$ are tested, and a triplet is added when the three bias labels form \{left, center, right\}.

This anchor iteration is the exact reason duplicate triplets were produced.
The same semantic group of three articles can be rediscovered when a different member of the same group is used as the anchor.
Because triplets are written in bias-slot form (left/center/right), rediscovering the same 3-article set can produce an exact duplicate row.
In practice, this created copies that appeared either twice or three times.
These exact duplicates were removed, reducing 6,714 rows to 4,849 unique triplets (saved in \texttt{data/unique\_triplets.json}).

To reduce long-tail over-representation from highly connected articles, we further applied a participation cap of 5 triplets per article.
This optimization step drops triplets while maximizing retained coverage under the cap constraint.
The capped dataset contains 2,631 triplets (saved in \texttt{data/capped\_triplets.json}), i.e., 2,218 triplets removed from the 4,849 unique-triplet set.

After capping, 3,485 unique articles participate in triplets, and the mean participation is
\[
\text{avg triplets per article} = \frac{3 \times 2631}{3485} = 2.26.
\]

The participation rate after capping is
\[
\frac{3485}{37554} \times 100 = 9.28\%.
\]

The fact that only about 9--10\% of raw articles participate is expected from the construction constraints in \texttt{data/sort\_data.py}:
(1) only top-$k$ nearest neighbors are considered ($k=10$),
(2) neighbors must pass a strict cosine similarity threshold $>0.8$,
and (3) a valid candidate must simultaneously satisfy the left-center-right label requirement.
Together, these filters favor high-confidence event matches and exclude most articles that do not have sufficiently similar cross-bias counterparts.

We apply the cap for the following reasons:
\begin{itemize}
\item Reduces over-representation from a few highly connected articles.
\item Makes training signal less skewed toward repeated anchors/events.
\item Usually improves robustness and fairness of comparisons across events/bias groups.
\item Gives cleaner interpretation when triplets are treated as event-level groupings.
\end{itemize}

\begin{table}[h]
\centering
\begin{tabular}{cc}
\hline
\textbf{Number of Articles} & \textbf{Participated in How Many Triplets} \\
\hline
1587 & 1 \\
741 & 2 \\
366 & 3 \\
229 & 4 \\
562 & 5 \\
\hline
\end{tabular}
\caption{Article participation distribution after capping at 5 triplets per article.}
\label{tab:capped_triplet_participation}
\end{table}

\subsection{Operational Valid-Triplet Subset (RST-Available Only)}

For model training without waiting for pending RST parses, we define an operational subset by excluding any triplet
that contains an article listed in \texttt{data/remaining\_articles.json}.
This yields a strictly RST-available triplet set.

Starting from 2,631 capped triplets, 883 triplets are excluded and 1,748 triplets remain valid
(stored in \texttt{data/valid\_triplets.json}).
The resulting valid subset contains 2,617 unique articles
(auxiliary list archived at \texttt{data/optional/extra/valid\_unique\_articles.json}).

This subset is used for immediate splitting, fact clustering, DFI construction, and SVM training.

\subsection{Leakage-safe Train/Val/Test Split}

We split the 1,748 valid triplets into train/validation/test with target ratio 80/10/10,
under a hard no-overlap constraint: no article ID may appear across different splits.
The split is solved at connected-component level over the triplet-document overlap graph.

The split is exactly feasible and yields:
\begin{itemize}
\item Train: 1,398 triplets
\item Validation: 174 triplets
\item Test: 176 triplets
\end{itemize}

The achieved proportions are 0.7998 / 0.0995 / 0.1007 (train/val/test),
and document overlap is zero across all split pairs.
Split metadata is stored in \texttt{data/valid\_triplet\_splits/meta.json}.

\section{EDU Filtering Pipeline}

The original RST output is rich but noisy: not every EDU is a standalone factual unit.
Since downstream fact alignment assumes that candidate EDUs correspond to factual propositions,
we introduced an explicit filtering stage in \texttt{modules/FactCluster.py} before SBERT encoding and clustering.
This reduces non-factual fragments (metadata, social-media markup, and short attribution shells)
that otherwise inflate clusters and distort prominence-based DFI computation.

\subsection{Pre-filter Audit (Before Filtering)}

Before implementing the filter, we audited the generated RST artifacts (3,005 documents)
and observed the following distributional properties:

\begin{table}[h]
\centering
\begin{tabular}{lr}
\hline
\textbf{Metric (Before Filtering)} & \textbf{Value} \\
\hline
Total EDUs & 353,410 \\
Average EDU length (tokens) & 8.41 \\
Median EDU length (tokens) & 7 \\
EDUs with $\leq 3$ tokens & 53,938 (15.26\%) \\
EDUs with $\leq 2$ tokens & 24,795 (7.02\%) \\
URL-like EDUs & 302 (0.09\%) \\
Punctuation-only EDUs & 560 (0.16\%) \\
Short attribution-like EDUs & 21,717 (6.14\%) \\
\hline
\end{tabular}
\caption{Global EDU audit on RST outputs before EDU candidate filtering.}
\label{tab:edu_prefilter_audit}
\end{table}

In practice, these artifacts included examples such as \textit{``Story highlights''},
\textit{``JUST WATCHED''}, \textit{``More:''}, social-media/link fragments, and short shells
like \textit{``he said''} that carry little standalone factual content.

\subsection{Implemented Filtering Logic}

Filtering is applied in \texttt{idify\_edus(...)} via
\texttt{\_is\_edu\_fact\_candidate(...)}.
The rules are configurable in \texttt{params.yaml} under \texttt{fact\_clustering.edu\_filter}.

\begin{itemize}
\item \textbf{Minimum lexical content}:
	EDUs with fewer than 3 tokens are dropped.
	This suppresses fragments that are unlikely to represent complete facts.

\item \textbf{Punctuation-only removal}:
	EDUs consisting only of punctuation/symbols are dropped,
	because they are parser artifacts rather than semantic units.

\item \textbf{URL and social markup removal}:
	URL-like/text-markup patterns (e.g., \texttt{https://}, \texttt{pic.twitter}, HTML-style links)
	are removed to avoid clustering on formatting or source metadata.

\item \textbf{Boilerplate/meta phrase removal}:
	editorial wrappers such as \textit{``Story highlights''}, \textit{``Read more''}, and similar
	non-propositional cues are filtered out.

\item \textbf{Short attribution-shell removal}:
	short EDUs dominated by attribution verbs (e.g., \textit{said/told/according to}) are removed,
	since they often encode discourse framing without independent factual content.
\end{itemize}

\subsection{Cluster-level Constraints After Filtering}

Beyond per-EDU filtering, we also tightened cluster retention:

\begin{itemize}
\item \textbf{Cross-bias minimum}:
	retained clusters must include center and at least one non-center side
	(\texttt{fact\_clustering.min\_biases\_per\_cluster = 2}).
	This allows center-left, center-right, and center-left-right clusters,
	which directly supports left-vs-center and right-vs-center DFI construction.

\item \textbf{Center-only singleton suppression}:
	singleton retention is disabled under cross-bias constraints,
	preventing isolated center-only fragments from entering the fact set.

\item \textbf{No representative-EDU downsampling}:
	for each retained cluster, all validated EDUs are kept.
	We do not collapse to one EDU per bias side.
\end{itemize}

Overall, the filtering pipeline improves precision of fact candidates,
aligns EDU selection with the intended notion of ``matchable fact'',
and reduces noise propagation into prominence scoring and DFI vectors.

\subsection{Operational Methodology: Fresh Re-clustering and Training (GPU)}

After creating the valid subset and leakage-safe split files, we execute the downstream pipeline
from repository root in four command stages (environment creation and dependency installation are omitted here):

\begin{enumerate}
\item \textbf{Fresh reclustering on valid triplets}
\begin{verbatim}
python run_fact_clustering.py \
	--triplets data/valid_triplets.json \
	--out data/valid_cluster_results_recluster_gpu.json \
	--meta data/valid_cluster_results_recluster_gpu_meta.json
\end{verbatim}

\item \textbf{Fact construction from fresh clusters}
\begin{verbatim}
python data/build_facts_from_clusters.py \
	--clusters data/valid_cluster_results_recluster_gpu.json \
	--out data/valid_facts_results_recluster_gpu.json \
	--meta data/valid_facts_results_recluster_gpu_meta.json
\end{verbatim}

\item \textbf{DFI split construction using predefined leakage-safe triplet splits}
\begin{verbatim}
python build_dfi_from_splits.py \
	--facts data/valid_facts_results_recluster_gpu.json \
	--split-dir data/valid_triplet_splits \
	--out-dir data/valid_dfi_splits_recluster_gpu
\end{verbatim}

\item \textbf{SVM training and evaluation}
\begin{verbatim}
python train_svm_from_dfi_splits.py \
	--split-dir data/valid_dfi_splits_recluster_gpu \
	--out data/valid_svm_metrics_recluster_gpu.json
\end{verbatim}
\end{enumerate}

Operational notes used in this run:
\begin{itemize}
\item Fresh clustering enforces center-required retention with \(\texttt{min\_biases\_per\_cluster}=2\)
			(center+left, center+right, or center+left+right).
\item All validated EDUs are retained in kept clusters (no representative-EDU downsampling).
\item Fresh clustering output files are guarded against accidental overwrite unless \texttt{--overwrite} is supplied.
\end{itemize}

\subsection{Fresh Valid-Subset Fact Clustering Results}

From 1,748 valid triplets, fresh reclustering produces:
\begin{itemize}
\item Triplets with non-empty retained clusters: 1,730
\item Empty/dropped triplets: 18
\item Errors: 0
\item Total retained clusters: 16,170
\item Total clustered EDU mentions: 47,668
\end{itemize}

EDU filtering statistics during clustering:
\begin{itemize}
\item Input EDU mentions: 630,516
\item Kept EDU mentions: 548,960
\item Dropped EDU mentions: 81,556
\item Drop ratio: 0.1293
\end{itemize}

Outputs:
\texttt{data/valid\_cluster\_results\_recluster\_gpu.json} and
\texttt{data/valid\_cluster\_results\_recluster\_gpu\_meta.json}.

\subsection{Fact Construction from Fresh Clusters}

Fact construction is run directly on the fresh cluster artifact.
All 1,730 cluster rows are converted into fact rows (errors: 0),
covering 16,170 clusters.

Outputs:
\begin{itemize}
\item Facts: \texttt{data/valid\_facts\_results\_recluster\_gpu.json}
\item Metadata: \texttt{data/valid\_facts\_results\_recluster\_gpu\_meta.json}
\end{itemize}

\subsection{DFI Split Construction from Predefined Triplet Splits}

DFI vectors are generated from
\texttt{data/valid\_facts\_results\_recluster\_gpu.json}
and assigned using
\texttt{data/valid\_triplet\_splits/train.json},
\texttt{val.json}, and \texttt{test.json}.

All 1,730 fact rows are matched to one split partition (unmatched: 0, errors: 0), yielding:
\begin{itemize}
\item Train DFI rows: 1,383
\item Validation DFI rows: 171
\item Test DFI rows: 176
\end{itemize}

DFI outputs are stored in
\texttt{data/valid\_dfi\_splits\_recluster\_gpu/}
with metadata in
\texttt{data/valid\_dfi\_splits\_recluster\_gpu/meta.json}.

\subsection{How DFI Vectors Are Built}

Each triplet uses one shared set of retained clusters.
For each cluster $c$, side-level prominence is computed for left, center, and right.
If multiple EDUs from the same side occur in a cluster, the side score is the maximum over that side.
If a side is absent in cluster $c$, its score is set to zero (omission case).

Using side scores $W^{(c)}_{left}$, $W^{(c)}_{center}$, and $W^{(c)}_{right}$,
the per-cluster DFI deltas are:
\[
\Delta^{(c)}_{left}=W^{(c)}_{left}-W^{(c)}_{center},\qquad
\Delta^{(c)}_{right}=W^{(c)}_{right}-W^{(c)}_{center}.
\]

The triplet-level vectors are then formed by concatenating these deltas across all $K$ retained clusters:
\[
\mathrm{DFI}_{left}=[\Delta^{(1)}_{left},\ldots,\Delta^{(K)}_{left}],\qquad
\mathrm{DFI}_{right}=[\Delta^{(1)}_{right},\ldots,\Delta^{(K)}_{right}].
\]
Therefore, both vectors always have the same length:
\[
|\mathrm{DFI}_{left}|=|\mathrm{DFI}_{right}|=K.
\]

For binary training, each triplet contributes exactly two samples:
\begin{itemize}
\item left-vs-center with label 0 (negative)
\item right-vs-center with label 1 (positive)
\end{itemize}
Hence for $N$ triplets, the class counts are exactly balanced:
\[
\#(\text{label }0)=N,\qquad \#(\text{label }1)=N.
\]

Before SVM training, all DFI vectors are converted to a fixed input dimension by right-padding with zeros
to the maximum training-vector length (and truncating longer vectors if needed at evaluation time).

\subsection{SVM Protocol and Evaluation (Fresh Re-clustering)}

The SVM stage uses binary side-comparison targets:
left-vs-center examples are labeled 0 and right-vs-center examples are labeled 1,
so each triplet contributes exactly two examples.

With 1,383/171/176 triplets in train/val/test DFI files, the effective binary example counts are:
\begin{itemize}
\item Train: 2,766
\item Validation: 342
\item Test: 352
\end{itemize}

Positive/negative distribution is exactly balanced in each split
(negative = left-vs-center, label 0; positive = right-vs-center, label 1):
\begin{table}[h]
\centering
\begin{tabular}{lccc}
\hline
\textbf{Split} & \textbf{Negative (label 0)} & \textbf{Positive (label 1)} & \textbf{Ratio} \\
\hline
Train & 1,383 & 1,383 & 50\% / 50\% \\
Validation & 171 & 171 & 50\% / 50\% \\
Test & 176 & 176 & 50\% / 50\% \\
\hline
\end{tabular}
\caption{Exact class balance induced by two-sample construction per triplet.}
\label{tab:binary_class_balance}
\end{table}

Using RBF SVM (\(C=10\), \(\gamma=0.1\), degree 3), raw padded DFI features
(input dimension 69), we obtain:
\begin{itemize}
\item Validation accuracy: 0.6988, macro-F1: 0.6987
\item Test accuracy: 0.6705, macro-F1: 0.6702
\end{itemize}

Detailed metrics are stored in
\texttt{data/valid\_svm\_metrics\_recluster\_gpu.json}.

\subsection{Coverage-controlled Structural Ablation (Size-3 Clusters Only)}

To isolate structural weighting effects from cross-bias coverage/omission cues,
we run a stricter ablation where DFI is generated only from clusters that are
\textit{exactly size 3} and contain all three bias sides (left, center, right).
This removes center+left / center+right omission patterns from the feature construction.

Methodology for this control setting:
\begin{itemize}
\item Keep only exact size-3 tri-side clusters per triplet.
\item Drop triplets that have no retained cluster after this filtering.
\item Rebuild DFI vectors and rerun the same four experiments:
baseline, without $s$, without $d$, and without both.
\item Use the same SVM configuration (RBF, $C=10$, $\gamma=0.1$, degree 3).
\end{itemize}

After filtering, the usable triplet counts become:
\begin{itemize}
\item Train: 489 (from 1,383)
\item Validation: 84 (from 171)
\item Test: 72 (from 176)
\end{itemize}

Since each triplet still contributes two binary examples, the corresponding sample counts are:
Train 978, Validation 168, Test 144 (still exactly 50/50 label balance).

\begin{table}[h]
\centering
\begin{tabular}{lcccc}
\hline
\textbf{Experiment} & \textbf{Val Acc.} & \textbf{Val Macro-F1} & \textbf{Test Acc.} & \textbf{Test Macro-F1} \\
\hline
Baseline & 0.5238 & 0.5121 & 0.5347 & 0.5309 \\
Without $s$ & 0.5119 & 0.5049 & 0.5625 & 0.5615 \\
Without $d$ & 0.5179 & 0.4841 & 0.5139 & 0.4852 \\
Without both ($d,s$) & 0.5000 & 0.3333 & 0.5000 & 0.3333 \\
\hline
\end{tabular}
\caption{Size-3 tri-side ablation results (coverage-controlled setting).}
\label{tab:size3_structural_ablation_results}
\end{table}

Interpretation in the coverage-controlled setting:
\begin{itemize}
\item The large gain from ``without both'' seen in the unrestricted setting disappears.
\item ``Without both'' collapses to near-chance accuracy with macro-F1 $\approx 0.333$, indicating loss of discriminative structure signal.
\item Compared with baseline, removing depth ($d$) hurts more than removing satellite term ($s$), suggesting depth contributes useful information in this stricter setup.
\item The slight test improvement for ``without $s$'' suggests satellite-edge weighting may be noisier than depth under these constraints.
\item Absolute scores are lower than the unrestricted setting because the exact size-3 filter substantially reduces available training/evaluation data.
\end{itemize}

\subsection{Expected Fresh-Run Artifacts}

The complete fresh-run artifact set is:
\begin{itemize}
\item \texttt{data/valid\_cluster\_results\_recluster\_gpu.json}
\item \texttt{data/valid\_cluster\_results\_recluster\_gpu\_meta.json}
\item \texttt{data/valid\_facts\_results\_recluster\_gpu.json}
\item \texttt{data/valid\_facts\_results\_recluster\_gpu\_meta.json}
\item \texttt{data/valid\_dfi\_splits\_recluster\_gpu/train.json}
\item \texttt{data/valid\_dfi\_splits\_recluster\_gpu/val.json}
\item \texttt{data/valid\_dfi\_splits\_recluster\_gpu/test.json}
\item \texttt{data/valid\_dfi\_splits\_recluster\_gpu/meta.json}
\item \texttt{data/valid\_svm\_metrics\_recluster\_gpu.json}
\end{itemize}

\section{Coverage-Only Baseline Experiment (Option A)}

The structural ablation results showed that removing both structural parameters ($\alpha=1$, $\gamma=1$)
achieved the highest test accuracy (75.28\%) in the unrestricted setting.
This raised a key question: is the model learning from RST structural signals at all,
or is it primarily exploiting fact coverage/omission patterns?

To answer this, we implemented a \textbf{pure coverage-only model} that encodes only the binary
presence or absence of facts per bias side, with no structural weighting whatsoever.

\subsection{Coverage-Only Feature Construction}

For each cluster $c$ in a triplet, we compute binary presence indicators:
\[
h^{(c)}_{side} = \begin{cases} 1 & \text{if side has at least one EDU in cluster } c \\ 0 & \text{otherwise} \end{cases}
\]
where $side \in \{left, center, right\}$.

The coverage delta for each cluster mirrors the DFI delta structure:
\[
\delta^{(c)}_{left} = h^{(c)}_{left} - h^{(c)}_{center}, \qquad
\delta^{(c)}_{right} = h^{(c)}_{right} - h^{(c)}_{center}
\]

These deltas take values in $\{-1, 0, +1\}$:
\begin{itemize}
\item $-1$: side omits a fact that center covers
\item $0$: both cover or both omit
\item $+1$: side covers a fact that center omits
\end{itemize}

The triplet-level coverage vectors are formed by concatenating deltas across all $K$ clusters:
\[
\mathrm{Cov}_{left} = [\delta^{(1)}_{left}, \ldots, \delta^{(K)}_{left}], \qquad
\mathrm{Cov}_{right} = [\delta^{(1)}_{right}, \ldots, \delta^{(K)}_{right}]
\]

This encoding is mathematically equivalent to the DFI construction with $W(d,s) = 1$ for present EDUs
and $W(d,s) = 0$ for absent EDUs---i.e., the ``without both'' condition.

\subsection{Coverage-Only Results}

We trained an SVM with identical configuration (RBF, $C=10$, $\gamma=0.1$) on coverage-only features.
The results confirm that coverage patterns fully explain the performance of the ``without both'' model:

\begin{table}[h]
\centering
\begin{tabular}{lcccc}
\hline
\textbf{Model} & \textbf{Val Acc.} & \textbf{Val F1} & \textbf{Test Acc.} & \textbf{Test F1} \\
\hline
Structural baseline ($\alpha=0.8$, $\gamma=0.5$) & 0.6988 & 0.6987 & 0.6705 & 0.6702 \\
Without $s$ ($\gamma=1$) & 0.6813 & 0.6803 & 0.7017 & 0.7008 \\
Without $d$ ($\alpha=1$) & 0.7135 & 0.7134 & 0.6591 & 0.6591 \\
Without both ($\alpha=1$, $\gamma=1$) & 0.7865 & 0.7845 & 0.7528 & 0.7504 \\
\textbf{Coverage-only (binary delta)} & \textbf{0.7865} & \textbf{0.7845} & \textbf{0.7528} & \textbf{0.7504} \\
\hline
\end{tabular}
\caption{Comparison of structural DFI models with the coverage-only baseline.}
\label{tab:coverage_comparison}
\end{table}

\subsection{Interpretation of Coverage-Only Results}

The coverage-only model achieves \textbf{identical performance} to the ``without both'' structural model
(test accuracy 75.28\%, macro-F1 75.04\%). This confirms several important findings:

\begin{enumerate}
\item \textbf{The ``without both'' model learns purely from omission patterns.}
      When $\alpha=\gamma=1$, the prominence score $W(d,s) = 1^{d+1} \cdot 1^s = 1$ for all present EDUs,
      reducing the DFI delta to a binary coverage signal.

\item \textbf{RST structural features provide no additional predictive value} when coverage information is available.
      The structural baseline (67.05\%) actually \textit{underperforms} the coverage-only model (75.28\%)
      by 8.23 percentage points.

\item \textbf{Structural weighting introduces noise rather than signal.}
      The depth and satellite parameters modulate the coverage signal in ways that reduce,
      rather than enhance, discriminative power.

\item \textbf{Fact coverage asymmetry is the dominant signal for bias classification.}
      Different bias orientations systematically cover or omit different facts,
      and this coverage pattern alone achieves strong classification performance.
\end{enumerate}

This result motivates a reframing of the research contribution:
rather than ``RST structure predicts bias,'' the finding is that
\textbf{``fact omission patterns across left/center/right sources predict bias with 75\% accuracy.''}

Outputs are stored in \texttt{omission-based/results/coverage\_svm\_results.json}
and \texttt{omission-based/results/coverage\_svm\_model.pkl}.

\section{Strengthening Structural Signals (Option B)}

Given that the original prominence formula underperforms the coverage-only baseline,
we explored whether alternative structural feature encodings could capture predictive
information that the original $W(d,s) = \alpha^{d+1} \gamma^s$ formula misses.

Three approaches were systematically evaluated:
\begin{enumerate}
\item \textbf{Aggregate structural features}: Mean/variance of depth, nuclearity distributions
\item \textbf{Richer RST features}: Nuclearity role encoding per cluster
\item \textbf{Alternative prominence formulas}: Exponential, inverse, logarithmic variants
\end{enumerate}

\subsection{Approach 1: Aggregate Structural Features}

Instead of computing a single max-prominence score per side, we aggregate multiple structural
statistics across all EDUs belonging to each side:

\begin{itemize}
\item \textbf{Depth statistics}: mean, variance of RST tree depth
\item \textbf{Nuclearity ratio}: fraction of EDUs with nucleus role (vs. satellite)
\item \textbf{Satellite edge statistics}: mean count of satellite edges to root
\end{itemize}

For each cluster, we compute these statistics per side and encode deltas against center.
Two variants were tested:
\begin{itemize}
\item \textbf{With coverage}: includes presence/absence indicator (retains coverage signal)
\item \textbf{Without coverage}: only includes clusters where all three sides are present (isolates structural signal)
\end{itemize}

\subsection{Approach 2: Nuclearity Distribution Features}

This approach focuses specifically on the nucleus/satellite distinction in RST theory.
For each cluster, we encode:
\begin{itemize}
\item Binary coverage indicator per side
\item Nucleus ratio (proportion of EDUs that are nuclei) per side
\end{itemize}

The intuition is that biased sources may systematically place their preferred facts in nucleus
positions (rhetorically prominent) while relegating opposing facts to satellite positions.

\subsection{Approach 3: Alternative Prominence Formulas}

We tested four alternative weighting functions:

\begin{enumerate}
\item \textbf{Exponential decay}: $W = e^{-\beta d}$ with $\beta=0.3$

This replaces the geometric decay $\alpha^{d+1}$ with exponential decay,
providing smoother depth penalty.

\item \textbf{Inverse satellite}: $W = \frac{1}{1 + \gamma s}$ with $\gamma=0.5$

This removes depth entirely and uses only satellite edge count with inverse penalty.

\item \textbf{Combined}: $W = \frac{e^{-\beta d}}{1 + \gamma s}$

Combines exponential depth decay with inverse satellite penalty.

\item \textbf{Logarithmic depth}: $W = \frac{1}{1 + \ln(1+d)}$

Uses logarithmic depth penalty, which compresses deep differences relative to shallow ones.
\end{enumerate}

\subsection{Option B Experimental Results}

All experiments used identical SVM configuration (RBF, $C=10$, $\gamma=0.1$) and train/val/test splits.

\begin{table}[h]
\centering
\small
\begin{tabular}{lcccccc}
\hline
\textbf{Experiment} & \textbf{Dim} & \textbf{Val Acc} & \textbf{Test Acc} & \textbf{Test F1} & \textbf{$\Delta$ Cov.} & \textbf{$\Delta$ Str.} \\
\hline
\multicolumn{7}{l}{\textit{Approach 1: Aggregate Features}} \\
\quad With coverage & 345 & 0.6754 & 0.6847 & 0.6679 & $-0.068$ & $+0.014$ \\
\quad No coverage (size-3) & 76 & 0.5175 & 0.5312 & 0.5289 & $-0.222$ & $-0.139$ \\
\hline
\multicolumn{7}{l}{\textit{Approach 2: Nuclearity Distribution}} \\
\quad Nuclearity dist. & 138 & 0.7368 & 0.7500 & 0.7499 & $-0.003$ & $+0.080$ \\
\hline
\multicolumn{7}{l}{\textit{Approach 3: Alternative Formulas}} \\
\quad $W = e^{-0.3d}$ & 69 & 0.6579 & 0.6705 & 0.6694 & $-0.082$ & $\pm0.000$ \\
\quad $W = 1/(1+0.5s)$ & 69 & 0.7749 & 0.7500 & 0.7500 & $-0.003$ & $+0.080$ \\
\quad $W = e^{-0.3d}/(1+0.5s)$ & 69 & 0.6901 & 0.6506 & 0.6491 & $-0.102$ & $-0.020$ \\
\quad $W = 1/(1+\ln(1+d))$ & 69 & 0.7749 & \textbf{0.7642} & \textbf{0.7641} & $+0.011$ & $+0.094$ \\
\quad Original $\alpha^{d+1}\gamma^s$ & 69 & 0.6988 & 0.6705 & 0.6702 & $-0.082$ & $\pm0.000$ \\
\hline
\multicolumn{7}{l}{\textit{Reference Baselines}} \\
\quad Coverage-only & 69 & 0.7865 & 0.7528 & 0.7504 & --- & $+0.082$ \\
\quad Structural baseline & 69 & 0.6988 & 0.6705 & 0.6702 & $-0.082$ & --- \\
\hline
\end{tabular}
\caption{Option B results: exploring strengthened structural features. $\Delta$ Cov. = delta vs coverage-only baseline; $\Delta$ Str. = delta vs original structural baseline.}
\label{tab:option_b_results}
\end{table}

\subsection{Interpretation of Option B Results}

\paragraph{Key Finding 1: Logarithmic depth achieves the best structural performance.}
The $W = 1/(1+\ln(1+d))$ formula achieves \textbf{76.42\% test accuracy}, which is:
\begin{itemize}
\item $+9.37\%$ above the original structural baseline (67.05\%)
\item $+1.14\%$ above the coverage-only baseline (75.28\%)
\end{itemize}
This is the \textit{only} experiment that outperforms coverage-only, suggesting that
logarithmic depth compression captures useful structural information that raw coverage misses.

\paragraph{Key Finding 2: Inverse satellite also performs well.}
The $W = 1/(1+0.5s)$ formula achieves \textbf{75.00\%} test accuracy, matching coverage-only.
Combined with the nuclearity distribution experiment (also 75.00\%), this suggests that
the satellite/nucleus distinction contains information comparable to coverage.

\paragraph{Key Finding 3: Aggregate features without coverage fail.}
When the coverage signal is removed (size-3 only clusters), aggregate structural features
achieve only \textbf{53.12\%}---near chance. This confirms that:
\begin{itemize}
\item Structural features alone (without coverage) have minimal predictive power
\item The gains from ``with coverage'' aggregate features come from the coverage signal, not structural statistics
\end{itemize}

\paragraph{Key Finding 4: Depth-only formulas underperform.}
Pure exponential depth decay ($e^{-0.3d}$) achieves only 67.05\%, identical to the original baseline.
This suggests depth alone is not the useful structural component---rather, it's the
\textit{compression} of depth differences (via logarithm) that helps.

\paragraph{Key Finding 5: Combining depth and satellite hurts.}
The combined formula $e^{-0.3d}/(1+0.5s)$ achieves only 65.06\%, \textit{worse} than either component alone.
This indicates that depth and satellite signals interact destructively when multiplied.

\subsection{Comparison Summary}

\begin{table}[h]
\centering
\begin{tabular}{lcc}
\hline
\textbf{Model} & \textbf{Test Accuracy} & \textbf{Relative to Coverage} \\
\hline
$W = 1/(1+\ln(1+d))$ (log depth) & \textbf{76.42\%} & $+1.14\%$ \\
Coverage-only & 75.28\% & baseline \\
$W = 1/(1+0.5s)$ (inverse sat.) & 75.00\% & $-0.28\%$ \\
Nuclearity distribution & 75.00\% & $-0.28\%$ \\
Original structural & 67.05\% & $-8.23\%$ \\
Aggregate (no coverage) & 53.12\% & $-22.16\%$ \\
\hline
\end{tabular}
\caption{Final comparison of all approaches.}
\label{tab:final_comparison}
\end{table}

\subsection{Conclusions from Option B}

\begin{enumerate}
\item \textbf{The logarithmic depth formula provides marginal improvement} over coverage-only (+1.14\%),
      suggesting that depth information can add value when appropriately compressed.

\item \textbf{The original prominence formula is suboptimal.}
      The geometric weighting $\alpha^{d+1}\gamma^s$ performs 9.37\% worse than the best alternative.

\item \textbf{Coverage remains the dominant signal.}
      Most structural variants cannot match coverage-only performance; those that do
      (inverse satellite, nuclearity) essentially encode coverage-correlated information.

\item \textbf{Pure structural features without coverage fail.}
      This definitively answers the research question: RST structural placement \textit{alone}
      does not predict bias; the predictive power comes from coverage asymmetries.
\end{enumerate}

Outputs are stored in \texttt{strengthen-str/results/strengthen\_structural\_results.json}
and trained models in \texttt{strengthen-str/results/models/}.

\section{Aggregate Vector Experiment}

A potential concern with the padded DFI approach is that DFI vectors from different triplets
have different facts/clusters, so the $i$-th entry in one DFI vector may correspond to a
completely different fact than the $i$-th entry in another DFI vector.
This raises the question: does the model learn meaningful patterns, or does it exploit
artifacts of the padding scheme?

\subsection{Motivation: Dimension Mismatch Problem}

In the standard DFI construction:
\begin{itemize}
\item Each triplet has $K$ retained clusters (variable across triplets)
\item DFI vectors have length $K$, padded/truncated to a fixed maximum dimension
\item The $i$-th dimension corresponds to cluster $i$ in that triplet's cluster set
\item Cluster identity is \textit{not} aligned across triplets---different triplets may have
      completely different facts in position $i$
\end{itemize}

The padding approach implicitly assumes that the SVM can learn from positional statistics
(e.g., ``position 5 tends to have high values for left-biased triplets''), even though
positions are not semantically aligned.

\subsection{Proposed Solution: Aggregate Feature Extraction}

To test whether aggregate statistics are sufficient, we extract fixed-length feature vectors
from each DFI that are semantically consistent across all triplets:

\begin{enumerate}
\item \textbf{Normalize} each raw DFI vector: $\hat{\Delta} = \frac{\Delta - \mu_\Delta}{\sigma_\Delta}$

\item \textbf{Extract aggregate statistics} (9 features per vector):
\begin{itemize}
\item Mean, standard deviation, minimum, maximum of deltas
\item Skewness (asymmetry of distribution)
\item Proportion of positive deltas (side $>$ center)
\item Proportion of zero deltas (equal coverage)
\item Proportion of negative deltas (center $>$ side)
\item Log number of clusters: $\ln(1 + K)$
\end{itemize}

\item \textbf{Train SVM} on these 9-dimensional fixed vectors instead of variable-length padded DFI
\end{enumerate}

\subsection{Aggregate Vector Experimental Results}

Six experiments were conducted with variations on coverage vs.\ structural base features
and normalization:

\begin{table}[h]
\centering
\small
\begin{tabular}{lccccc}
\hline
\textbf{Experiment} & \textbf{Dim} & \textbf{Val Acc} & \textbf{Test Acc} & \textbf{$\Delta$ Cov-Pad} & \textbf{$\Delta$ LogD-Pad} \\
\hline
agg\_log\_depth & 9 & 0.5409 & \textbf{0.6136} & $-0.139$ & $-0.151$ \\
agg\_combined & 18 & 0.5439 & 0.6080 & $-0.145$ & $-0.156$ \\
agg\_extended & 14 & 0.5322 & 0.5767 & $-0.176$ & $-0.188$ \\
agg\_coverage\_norm & 9 & 0.5585 & 0.5739 & $-0.179$ & $-0.190$ \\
agg\_log\_depth\_norm & 9 & 0.5263 & 0.5682 & $-0.185$ & $-0.196$ \\
agg\_coverage & 9 & 0.5936 & 0.5568 & $-0.196$ & $-0.207$ \\
\hline
\multicolumn{6}{l}{\textit{Reference baselines (padded DFI approach)}} \\
Coverage-only padded & 69 & 0.7865 & 0.7528 & --- & $-0.011$ \\
Log-depth padded & 69 & 0.7749 & 0.7642 & $+0.011$ & --- \\
\hline
\end{tabular}
\caption{Aggregate vector results. $\Delta$ Cov-Pad = delta vs.\ coverage-only padded; $\Delta$ LogD-Pad = delta vs.\ log-depth padded.}
\label{tab:aggregate_results}
\end{table}

\subsection{Interpretation: Negative Result}

The aggregate feature approach performs \textbf{substantially worse} than the padded DFI approach:

\begin{itemize}
\item Best aggregate model (agg\_log\_depth): 61.36\% test accuracy
\item Coverage-only padded: 75.28\% test accuracy
\item Gap: $-13.92$ percentage points
\end{itemize}

This negative result provides important insight:

\paragraph{Finding 1: Position-specific patterns matter.}
The large performance gap indicates that the padded DFI approach is \textit{not} merely
exploiting positional artifacts. The raw DFI values encode cluster-specific patterns
that are lost when aggregating to summary statistics.

\paragraph{Finding 2: Which facts are omitted matters more than how many.}
Aggregate features like ``proportion of positive deltas'' capture \textit{how much}
a side differs from center, but not \textit{which specific facts} differ.
The padded approach preserves cluster-level identity through consistent intra-triplet
ordering, allowing the model to learn patterns like ``when cluster $i$ is omitted,
this predicts bias $X$.''

\paragraph{Finding 3: The dimension mismatch is less problematic than expected.}
Although cross-triplet cluster alignment is imperfect, the SVM learns useful patterns
from the distribution of values within each vector. The zeros from padding encode
``absence of later clusters,'' which is itself informative about triplet complexity.

\paragraph{Finding 4: Normalization does not help.}
Normalizing DFI vectors before aggregation provides no benefit (sometimes hurts),
suggesting that the absolute scale of deltas carries information.

\subsection{Conclusions from Aggregate Vector Experiment}

\begin{enumerate}
\item \textbf{Aggregate statistics are insufficient} for bias classification.
      The 9-14 feature vectors lose too much information compared to raw padded DFI.

\item \textbf{Specific fact patterns are essential.}
      The model needs to know \textit{which} facts are omitted/emphasized, not just
      summary statistics of omission rates.

\item \textbf{The padded DFI approach is vindicated.}
      Despite theoretical concerns about dimension alignment, the empirical results
      show that padded DFI captures meaningful signal that aggregation destroys.

\item \textbf{Future work} could explore cluster-alignment strategies (e.g., matching
      clusters across triplets by semantic similarity) to enable more principled
      cross-triplet comparisons.
\end{enumerate}

Outputs are stored in \texttt{aggregate-vector/results/aggregate\_svm\_results.json}
and trained models in \texttt{aggregate-vector/results/models/}.

\section{Hybrid Approach: Separating Coverage and Structural Signals (Option C)}

Having established that coverage dominates and that the logarithmic depth formula provides
marginal improvement, we sought to explicitly quantify the independent contributions of
coverage and structural signals. Three approaches were implemented.

\subsection{Approach 1: Feature Group Ablation}

We trained separate models on four feature configurations:
\begin{enumerate}
\item \textbf{Coverage-only}: Binary presence/absence deltas per cluster
\item \textbf{Structural (all clusters)}: Log-depth prominence deltas, all clusters
\item \textbf{Structural (size-3 only)}: Log-depth deltas, only clusters with all 3 sides present (removes coverage signal)
\item \textbf{Combined}: Coverage and structural features interleaved per cluster (2 features per cluster)
\end{enumerate}

\begin{table}[h]
\centering
\begin{tabular}{lcccc}
\hline
\textbf{Configuration} & \textbf{Dim} & \textbf{Test Acc} & \textbf{Test F1} & \textbf{$\Delta$ vs Cov.} \\
\hline
Coverage-only & 69 & 75.28\% & 75.04\% & --- \\
Structural (all clusters) & 69 & 76.42\% & 76.41\% & $+1.14\%$ \\
Structural (size-3 only) & var & 51.42\% & 48.70\% & $-23.86\%$ \\
\textbf{Combined} & 138 & \textbf{77.27\%} & \textbf{77.27\%} & $+1.99\%$ \\
\hline
\end{tabular}
\caption{Feature group ablation results. Combined features achieve the best performance.}
\label{tab:hybrid_ablation}
\end{table}

\paragraph{Key Finding:} The combined model achieves \textbf{77.27\%} test accuracy---the highest
observed across all experiments. This is $+1.99\%$ above coverage-only and $+0.85\%$ above
structural (all clusters), demonstrating that coverage and structural features provide
complementary information.

\subsection{Approach 2: Two-Stage Classifier and Error Analysis}

To understand how structural features could complement coverage, we analyzed coverage model errors:

\begin{itemize}
\item Coverage model test errors: 87/352 (24.7\%)
\item Structural model could correct: 21/87 (24.1\%)
\item Both models wrong: 66/87 (75.9\%)
\end{itemize}

This reveals that while structural features can correct some coverage errors, the majority
of errors are \textit{shared}---cases where both coverage and structural signals fail.

We tested combination strategies that defer to structural predictions when coverage confidence is low:

\begin{table}[h]
\centering
\begin{tabular}{lccc}
\hline
\textbf{Strategy} & \textbf{Test Acc} & \textbf{$\Delta$ vs Cov.} & \textbf{Notes} \\
\hline
Coverage only & 75.28\% & --- & Baseline \\
Threshold $p < 0.60$ & 75.85\% & $+0.57\%$ & 34 uncertain samples \\
Weighted (0.8 cov, 0.2 str) & 76.42\% & $+1.14\%$ & Soft combination \\
Weighted (0.7 cov, 0.3 str) & 76.42\% & $+1.14\%$ & Optimal weighting \\
Weighted (0.5 cov, 0.5 str) & 75.85\% & $+0.57\%$ & Equal weights underperform \\
\hline
\end{tabular}
\caption{Two-stage combination strategies. Optimal weighting favors coverage (70--80\%).}
\label{tab:hybrid_twostage}
\end{table}

\paragraph{Key Finding:} The optimal combination weights coverage at 70--80\% and structural
at 20--30\%, confirming coverage remains the dominant signal while structural provides
modest complementary information.

\subsection{Approach 3: Stacking Ensemble}

We trained a meta-classifier (logistic regression) on the class probabilities from
independent coverage and structural models:

\begin{table}[h]
\centering
\begin{tabular}{lcc}
\hline
\textbf{Model} & \textbf{Test Acc} & \textbf{$\Delta$ vs Cov.} \\
\hline
Coverage base model & 75.28\% & --- \\
Structural base model & 76.42\% & $+1.14\%$ \\
\textbf{Stacked meta-model} & \textbf{76.99\%} & $+1.71\%$ \\
\hline
\end{tabular}
\caption{Stacking ensemble results.}
\label{tab:hybrid_stacking}
\end{table}

The meta-model coefficients reveal relative feature importance:
\begin{itemize}
\item Coverage probabilities: $-1.21$ (class 0), $+1.25$ (class 1)
\item Structural probabilities: $-1.41$ (class 0), $+1.45$ (class 1)
\end{itemize}

Interestingly, structural features have slightly higher magnitude coefficients,
suggesting they provide marginally more discriminative power when both signals are available.

\subsection{Conclusions from Hybrid Approach (Option C)}

\begin{enumerate}
\item \textbf{Combined features achieve the best performance.}
      Simple feature concatenation (coverage + log-depth structural) yields \textbf{77.27\%}---the
      highest accuracy across all experiments, confirming that structural information adds
      value beyond coverage.

\item \textbf{Coverage and structure are partially redundant.}
      75.9\% of errors are shared between models, indicating substantial overlap in
      what they capture.

\item \textbf{Coverage remains dominant.}
      Optimal combination weights favor coverage (70--80\%), and the stacking improvement
      over coverage-only is modest (+1.71\%).

\item \textbf{Pure structural without coverage fails.}
      When coverage signal is removed (size-3 clusters only), accuracy drops to 51.42\%
      (near chance), definitively confirming that structural placement alone cannot predict bias.

\item \textbf{Practical recommendation:}
      For best performance, use combined features (coverage + log-depth structural).
      The $+2\%$ improvement over coverage-only justifies the added complexity.
\end{enumerate}

Outputs are stored in \texttt{hybrid-approach/results/hybrid\_results.json}.

\section{Experimental Design Improvements (Option D)}

The original fixed train/val/test split may have high variance due to relatively small test set
size (176 triplets = 352 samples). To provide more robust evaluation and quantify uncertainty,
we implemented several experimental design improvements.

\subsection{K-Fold Cross-Validation}

We replaced the fixed 80/10/10 split with stratified 5-fold cross-validation across the full
dataset of 1,730 triplets (3,460 samples). This reduces variance from lucky/unlucky test set
compositions.

\begin{table}[h]
\centering
\begin{tabular}{lcccc}
\hline
\textbf{Feature Set} & \textbf{Mean Acc.} & \textbf{Std} & \textbf{Min} & \textbf{Max} \\
\hline
Coverage & 77.34\% & 0.85\% & 76.16\% & 78.76\% \\
Structural (log-depth) & 77.02\% & 0.90\% & 75.72\% & 77.75\% \\
Combined & 76.65\% & 1.35\% & 74.71\% & 78.47\% \\
\hline
\end{tabular}
\caption{5-fold cross-validation results. Standard deviations are relatively small ($<1.5\%$),
indicating stable performance across folds.}
\label{tab:cv_results}
\end{table}

\paragraph{Key Finding:} The cross-validation results show consistent performance across folds
with low variance (std $<1.5\%$). Interestingly, coverage-only achieves the highest mean accuracy
(77.34\%) in the CV setting, slightly outperforming structural and combined features. This differs
from the fixed-split results where combined features achieved 77.27\%, suggesting some sensitivity
to data partitioning.

\subsection{Bootstrap Confidence Intervals}

To quantify uncertainty in accuracy estimates, we computed bootstrap 95\% confidence intervals
by resampling test predictions 1,000 times with replacement:

\begin{table}[h]
\centering
\begin{tabular}{lcccc}
\hline
\textbf{Feature Set} & \textbf{Point Est.} & \textbf{95\% CI} & \textbf{CI Width} & \textbf{SE} \\
\hline
Coverage & 78.47\% & [75.43\%, 81.65\%] & 6.21\% & 1.58\% \\
Structural & 76.30\% & [73.27\%, 79.34\%] & 6.07\% & 1.61\% \\
Combined & 75.43\% & [72.40\%, 78.76\%] & 6.36\% & 1.62\% \\
\hline
\end{tabular}
\caption{Bootstrap 95\% confidence intervals. CI widths of $\approx6\%$ indicate moderate
uncertainty in point estimates.}
\label{tab:bootstrap_ci}
\end{table}

\paragraph{Key Finding:} The confidence intervals are moderately wide ($\approx6\%$), indicating
that point estimates should be interpreted with caution. The overlapping CIs suggest that
differences between feature sets may not be statistically significant.

\subsection{Paired Bootstrap Test: Combined vs Coverage}

To formally test whether combined features significantly outperform coverage-only, we performed
a paired bootstrap test:

\begin{itemize}
\item Observed difference: $-3.03\%$ (combined minus coverage)
\item 95\% CI for difference: $[-5.35\%, -0.72\%]$
\item One-sided p-value: 0.997
\item Significant at $\alpha=0.05$: \textbf{No}
\end{itemize}

\paragraph{Key Finding:} In this particular train/test split, the combined features actually
\textit{underperform} coverage-only by 3\%. The 95\% CI for the difference excludes zero in
the wrong direction, indicating that the combined improvement observed in the fixed split
may not be robust.

\subsection{Learning Curve Analysis}

To assess whether more training data would improve performance, we generated learning curves
by training on increasing fractions of the training set:

\begin{figure}[h]
\centering
\begin{tabular}{cc}
\textbf{Train Fraction} & \textbf{Coverage Acc.} \\
\hline
10\% & 75.29\% \\
30\% & 75.14\% \\
50\% & 76.30\% \\
70\% & 76.88\% \\
100\% & 78.47\% \\
\end{tabular}
\caption{Learning curve for coverage features. Accuracy increases slowly with training size.}
\label{tab:learning_curve}
\end{figure}

\paragraph{Key Finding:} Performance at 10\% of training data (276 samples) is already 75.29\%,
only 3\% below full training set performance. This suggests the model learns quickly and
approaches a plateau, indicating that the current dataset size is likely adequate for this task.

\subsection{Conclusions from Option D}

\begin{enumerate}
\item \textbf{Results are moderately stable.} Cross-validation standard deviations of $<1.5\%$
indicate reasonable stability across different data partitions.

\item \textbf{Uncertainty is substantial.} Bootstrap 95\% CIs of $\approx6\%$ width mean
point estimates should be reported with confidence bounds.

\item \textbf{Combined vs coverage difference is not robust.} The paired test shows the
combined feature improvement is not statistically significant and may reverse in different splits.

\item \textbf{Dataset size is adequate.} Learning curves show performance plateaus early,
suggesting diminishing returns from additional training data.

\item \textbf{Practical recommendation:} Report accuracy as mean $\pm$ std from cross-validation
or as point estimate with 95\% CI, rather than single fixed-split numbers.
\end{enumerate}

Outputs are stored in \texttt{experimental-design/results/crossval\_results.json}.

\section{Alternative Models (Option E)}

To determine whether SVM with RBF kernel is optimal for this task, we compared multiple
model families with different inductive biases.

\subsection{Model Comparison}

We tested five model families on all three feature sets (coverage, structural, combined):

\begin{table}[h]
\centering
\small
\begin{tabular}{lccc}
\hline
\textbf{Model} & \textbf{Coverage} & \textbf{Structural} & \textbf{Combined} \\
\hline
SVM (RBF) & 76.70\% & 77.56\% & 77.84\% \\
Logistic (L2) & 74.43\% & 73.58\% & 75.28\% \\
Logistic (L1) & 75.00\% & 74.15\% & 76.14\% \\
MLP & 79.26\% & 78.12\% & 76.70\% \\
\textbf{Random Forest} & \textbf{79.26\%} & \textbf{86.08\%} & \textbf{86.36\%} \\
\hline
\end{tabular}
\caption{Model comparison across feature sets. Random Forest substantially outperforms all other models.}
\label{tab:model_comparison}
\end{table}

\paragraph{Key Finding:} \textbf{Random Forest dramatically outperforms SVM}, achieving 86.36\%
test accuracy on combined features---8.52 percentage points higher than SVM (77.84\%).
This is the largest improvement observed across all experiments.

\subsection{Hyperparameter Tuning}

We performed grid search with 5-fold cross-validation to tune key models on combined features:

\begin{table}[h]
\centering
\begin{tabular}{lcccc}
\hline
\textbf{Model} & \textbf{Best Params} & \textbf{CV Acc.} & \textbf{Test Acc.} \\
\hline
Logistic (tuned) & C=1.0, penalty=L2 & 76.93\% & 75.28\% \\
\textbf{Random Forest (tuned)} & depth=None, n=50, split=5 & 85.20\% & \textbf{87.78\%} \\
\hline
\end{tabular}
\caption{Hyperparameter tuning results on combined features.}
\label{tab:tuning_results}
\end{table}

\paragraph{Key Finding:} Tuned Random Forest achieves \textbf{87.78\% test accuracy}---the
highest performance observed in any experiment. The best configuration uses unlimited depth
with 50 trees and minimum 5 samples per split.

\subsection{Feature Importance Analysis}

Random Forest provides built-in feature importance via Gini impurity reduction.
Analysis of coverage features shows:

\begin{table}[h]
\centering
\begin{tabular}{cc}
\hline
\textbf{Feature Index} & \textbf{Importance} \\
\hline
0 & 0.3939 \\
1 & 0.1662 \\
2 & 0.0817 \\
3 & 0.0463 \\
4--9 & $<0.04$ each \\
\hline
\end{tabular}
\caption{Top feature importances for coverage features. Feature 0 dominates.}
\label{tab:feature_importance}
\end{table}

\paragraph{Key Finding:} The top 3 features account for $>64\%$ of total importance, suggesting
that a small number of cluster positions are highly predictive. Feature 0 alone contributes
39.4\% of the total importance, indicating one particular fact omission pattern is dominant.

\subsection{Interpretation of Random Forest Superiority}

Several factors may explain why Random Forest outperforms SVM:

\begin{enumerate}
\item \textbf{Non-linear interactions:} Random Forest can capture complex feature interactions
through tree splits, while SVM with RBF kernel may require careful hyperparameter tuning
to model similar interactions.

\item \textbf{Heterogeneous features:} The combined feature set mixes binary coverage indicators
with continuous structural scores. Tree-based models handle mixed feature types naturally,
while kernel methods may struggle with different scales.

\item \textbf{Sparse signal:} Most DFI values are zero (from padding). Decision trees can
efficiently ignore zero-valued features, while SVM must process all dimensions.

\item \textbf{Robustness to outliers:} Random Forest's ensemble averaging reduces sensitivity
to outliers and noise, potentially helping with the variable-length DFI vectors.
\end{enumerate}

\subsection{Conclusions from Option E}

\begin{enumerate}
\item \textbf{Model choice matters substantially.} Random Forest (87.78\%) outperforms
SVM (77.84\%) by nearly 10 percentage points---larger than any feature engineering improvement.

\item \textbf{SVM was suboptimal for this task.} The RBF SVM, while reasonable, significantly
underperforms tree-based methods on this sparse, heterogeneous feature space.

\item \textbf{Simple models underperform.} Logistic regression achieves only 75--76\%,
confirming that the classification boundary is non-linear.

\item \textbf{Combined features remain best with Random Forest.} The benefit of combining
coverage and structural features (+0.28\% over structural-only) persists with Random Forest.

\item \textbf{Practical recommendation:} Use Random Forest with unlimited depth, 50--100 trees,
and combined (coverage + log-depth structural) features for best performance.
\end{enumerate}

Outputs are stored in \texttt{alternative-models/results/alternative\_models\_results.json}.

\section{Summary of All Experimental Results}

\begin{table}[h]
\centering
\small
\begin{tabular}{llcc}
\hline
\textbf{Option} & \textbf{Best Configuration} & \textbf{Test Acc.} & \textbf{Key Finding} \\
\hline
Baseline & SVM + original $\alpha^{d+1}\gamma^s$ & 67.05\% & Original formula suboptimal \\
A & Coverage-only + SVM & 75.28\% & Omission patterns dominate \\
B & Log-depth structural + SVM & 76.42\% & +1.14\% over coverage \\
Aggregate & Summary stats + SVM & 61.36\% & Specific facts matter \\
C & Combined + SVM & 77.27\% & +1.99\% over coverage \\
D & Coverage + SVM (5-fold CV) & 77.34\% & Results are stable \\
\textbf{E} & \textbf{Combined + RF (tuned)} & \textbf{87.78\%} & \textbf{+10\% over SVM} \\
\hline
\end{tabular}
\caption{Summary of all experimental results. Random Forest with combined features achieves
the highest accuracy.}
\label{tab:full_summary}
\end{table}

\subsection{Key Takeaways}

\begin{enumerate}
\item \textbf{Fact omission is the dominant signal.} Coverage/omission patterns predict bias
better than RST structural placement across all experiments.

\item \textbf{Structural features provide modest improvement.} Log-depth prominence adds
$\approx1-2\%$ when combined with coverage, but fails alone without coverage signal.

\item \textbf{Model choice matters more than feature engineering.} Switching from SVM to
Random Forest improves accuracy by $\approx10\%$---far larger than any feature modification.

\item \textbf{The original prominence formula is suboptimal.} The logarithmic depth formula
$W = 1/(1+\ln(1+d))$ substantially outperforms the original $W = \alpha^{d+1}\gamma^s$.

\item \textbf{Aggregate statistics lose critical information.} The specific facts that are
omitted matter more than summary statistics of omission rates.

\item \textbf{Best practice:} Use Random Forest with combined features (coverage + log-depth
structural) for highest accuracy (87.78\%), or coverage-only for interpretability (79.26\% with RF).
\end{enumerate}

\section{Appendix: Feature Representation Clarifications}

This section clarifies the feature representations used across experiments, addressing common
questions about the relationship between different feature encodings.

\subsection{What Features Were Used for Training Random Forest (Option E)?}

Random Forest in Option E used \textbf{DFI-style padded vectors} (per-cluster deltas),
\textit{not} aggregate vectors. The features are computed per-cluster as deltas (side minus center),
then padded to a fixed length:

\begin{itemize}
\item \textbf{Coverage}: 69 dimensions (binary presence deltas)
\item \textbf{Structural}: 69 dimensions (log-depth prominence deltas)
\item \textbf{Combined}: 138 dimensions (both interleaved)
\end{itemize}

The aggregate vector experiment was a \textit{separate} negative result showing that summary
statistics (9--14 dimensions) lose too much information compared to the full padded DFI approach.

\subsection{Differences Between Coverage, Structural, and Combined Feature Sets}

\begin{table}[h]
\centering
\begin{tabular}{lll}
\hline
\textbf{Feature Set} & \textbf{Per-Cluster Value} & \textbf{Formula} \\
\hline
Coverage & Binary presence delta & $h_{side} - h_{center} \in \{-1, 0, +1\}$ \\
Structural & Prominence delta & $W_{side} - W_{center}$ where $W = \frac{1}{1+\ln(1+d)}$ \\
Combined & Both values & $[\text{coverage}_c, \text{structural}_c]$ per cluster $c$ \\
\hline
\end{tabular}
\caption{Feature set definitions. Coverage encodes binary presence/absence; structural encodes
log-depth weighted prominence; combined interleaves both per cluster.}
\label{tab:feature_definitions}
\end{table}

\subsection{Consistency of Feature Sets Across Experiments}

\textbf{Yes}---Options A, B, C, D, and E all use the same \textbf{per-cluster delta encoding
with zero-padding} to a fixed maximum length. The only variations are:

\begin{enumerate}
\item \textbf{Which delta type}: coverage (binary) vs structural (continuous) vs combined (both)
\item \textbf{Which weighting formula for structural}: original $\alpha^{d+1}\gamma^s$ vs
      log-depth $1/(1+\ln(1+d))$ vs inverse-satellite $1/(1+\gamma s)$
\end{enumerate}

The \textbf{aggregate vector experiment was separate}---it used summary statistics (mean, std,
skewness, etc.) instead of per-cluster vectors. This approach performed substantially worse
(61\% vs 75\%+), confirming that specific cluster-level patterns are essential.

\subsection{DFI Vectors vs Aggregate Vectors}

\begin{table}[h]
\centering
\begin{tabular}{lccl}
\hline
\textbf{Experiment} & \textbf{Encoding} & \textbf{Dim} & \textbf{Best Test Acc.} \\
\hline
Options A--E & DFI (per-cluster, padded) & 69--138 & 87.78\% (RF, combined) \\
Aggregate Vector & Summary statistics & 9--14 & 61.36\% \\
\hline
\end{tabular}
\caption{Comparison of DFI vs aggregate encoding approaches.}
\label{tab:dfi_vs_aggregate}
\end{table}

The key insight is that \textbf{which specific facts are omitted} matters far more than
aggregate statistics about omission rates. The per-cluster DFI encoding preserves this
information through consistent intra-triplet ordering, while aggregation destroys it.

\section{What Patterns Does the Model Learn?}

A natural question arises: \textit{what is the model actually learning?} Each triplet has
different fact clusters, so the $i$-th dimension in one DFI vector corresponds to a completely
different fact than the $i$-th dimension in another vector. Yet Random Forest achieves 87.78\%
accuracy, implying that meaningful cross-triplet patterns exist. This section investigates
what those patterns are.

\subsection{The Paradox of Unaligned Dimensions}

In the padded DFI approach:
\begin{itemize}
\item Each triplet has $K$ retained clusters (variable, mean $\approx 9$, max = 69)
\item DFI vectors are padded with zeros to a fixed length of 69
\item Cluster $i$ in triplet A may be about ``economic policy'' while cluster $i$ in triplet B
      is about ``immigration''---there is no semantic alignment
\end{itemize}

Despite this apparent misalignment, the model learns \textit{something} that generalizes.
The key insight is that the model does not learn ``cluster 3 predicts right-wing bias.''
Instead, it learns \textbf{distributional patterns} that are consistent across triplets.

\subsection{Pattern 1: Asymmetric Omission Rates}

The most robust pattern is that \textbf{right-leaning sources omit more facts that center covers}
compared to left-leaning sources:

\begin{table}[h]
\centering
\begin{tabular}{lcc}
\hline
\textbf{Sample Type} & \textbf{Mean Omissions} & \textbf{Std} \\
\hline
Left (label 0): left omits, center covers & 3.52 & 4.68 \\
Right (label 1): right omits, center covers & 3.97 & 4.31 \\
\hline
\end{tabular}
\caption{Omission statistics across training triplets. Right sources omit $\approx 13\%$ more
facts than left sources.}
\label{tab:omission_rates}
\end{table}

This asymmetry means that right-vs-center samples have systematically more negative deltas
(more $-1$ values in coverage vectors) than left-vs-center samples. The model can exploit
this aggregate pattern without knowing what specific facts are involved.

\subsection{Pattern 2: Position-Dependent Delta Distributions}

Analysis of the training data reveals that early positions (clusters 0--5) show stronger
systematic differences between left and right samples:

\begin{table}[h]
\centering
\small
\begin{tabular}{ccccc}
\hline
\textbf{Position} & \textbf{Left Mean} & \textbf{Right Mean} & \textbf{Difference} & \textbf{Non-zero Rate} \\
\hline
0 & $-0.069$ & $-0.566$ & $-0.497$ & 95--98\% \\
1 & $-0.159$ & $-0.504$ & $-0.345$ & 93--94\% \\
2 & $-0.222$ & $-0.452$ & $-0.230$ & 87--88\% \\
3 & $-0.244$ & $-0.422$ & $-0.177$ & 79\% \\
4 & $-0.239$ & $-0.350$ & $-0.111$ & 70\% \\
5 & $-0.238$ & $-0.273$ & $-0.035$ & 62\% \\
\hline
\end{tabular}
\caption{Coverage delta means by position for left (label 0) vs right (label 1) samples.
Early positions show larger systematic differences.}
\label{tab:position_deltas}
\end{table}

The key observations are:
\begin{enumerate}
\item \textbf{Both sides show negative means}: Both left and right sources omit facts that
      center covers, but right omits more.
\item \textbf{Differences decrease with position}: Early clusters (position 0--2) show the
      largest differences between labels; later positions converge.
\item \textbf{Non-zero rates decrease}: Later positions are mostly padding zeros, providing
      less signal.
\end{enumerate}

\subsection{Pattern 3: Vector Length as Implicit Feature}

The number of non-zero entries encodes triplet complexity. Triplets with more shared facts
(larger clusters) have longer DFI vectors. This implicit feature may correlate with bias
patterns---for example, highly contentious events may generate more coverage from all sides.

\subsection{Why Random Forest Excels}

Random Forest's architecture is particularly suited to exploit these patterns:

\begin{enumerate}
\item \textbf{Automatic feature selection}: Trees can learn to focus on early positions
      (where signal is strongest) and ignore later positions (mostly zeros). The feature
      importance analysis confirms this: position 0 accounts for 39.4\% of total importance.

\item \textbf{Threshold-based splits}: Decision trees naturally handle the discrete coverage
      deltas $\{-1, 0, +1\}$ through threshold splits, learning rules like
      ``if position 0 $< -0.3$ then predict right.''

\item \textbf{Robustness to sparse features}: Most DFI values are zero (padding). Trees
      efficiently ignore uninformative zeros, while kernel methods (SVM) must process all
      dimensions.

\item \textbf{Non-linear interactions}: Random Forest captures interactions like
      ``position 0 is negative AND position 1 is negative $\Rightarrow$ right,'' which
      linear models cannot represent.
\end{enumerate}

\subsection{What the Model Does NOT Learn}

It is important to clarify what the model is \textit{not} learning:

\begin{itemize}
\item \textbf{Not semantic fact content}: The model has no knowledge of what facts are
      about (immigration, economy, etc.). It only sees numerical deltas.

\item \textbf{Not cross-triplet fact alignment}: The model does not know that cluster 3
      in triplet A is semantically similar to cluster 7 in triplet B.

\item \textbf{Not RST structural placement}: The size-3 ablation showed that structural
      signals alone (without coverage) achieve only 51\% accuracy---near chance.
\end{itemize}

\subsection{Summary: The Learned Pattern}

The model learns a \textbf{distributional fingerprint} of bias-specific omission behavior:

\begin{quote}
\textit{Right-leaning sources systematically omit more facts that centrist sources cover,
and this omission pattern is particularly pronounced in the first few (most salient) facts
of each event.}
\end{quote}

This pattern is robust because:
\begin{enumerate}
\item It is a population-level tendency, not tied to specific facts
\item It exploits the inherent ordering of clusters within each triplet (larger/more important
      clusters tend to appear first)
\item Random Forest can efficiently learn the position-dependent thresholds that capture
      this asymmetry
\end{enumerate}

The 87.78\% accuracy reflects the model's ability to exploit these distributional regularities,
even without explicit semantic alignment across triplets.
